\documentclass[conference]{IEEEtran}
\IEEEoverridecommandlockouts

\usepackage{cite}
\usepackage{amsmath,amssymb,amsfonts}
\usepackage{algorithm}
\usepackage{algpseudocode}
\usepackage{graphicx}
\usepackage{textcomp}
\usepackage{xcolor}
\usepackage{booktabs}
\usepackage{array}
\usepackage{url}
\usepackage{hyperref}
\usepackage{balance}
\usepackage{multirow}

\hypersetup{
  colorlinks=true,
  linkcolor=blue,
  citecolor=blue,
  urlcolor=blue
}

\def\BibTeX{{\rm B\kern-.05em{\sc i\kern-.025em b}\kern-.08em
    T\kern-.1667em\lower.7ex\hbox{E}\kern-.125emX}}

\begin{document}

\title{Explainable Anomaly Detection in IoT System Logs\\
Using XGBoost, Random Forest, SMOTE, and SHAP}

\author{
  \IEEEauthorblockN{Faisal Ahamed Syed}
  \IEEEauthorblockA{
    \textit{Department of Computer Information Systems and Cybersecurity}\\
    \textit{Auburn University at Montgomery}\\
    Montgomery, AL, USA\\
    fsyed8@aum.edu
  }
}

\maketitle

% ────────────────────────────────────────────────────────────────────────
\begin{abstract}
Anomaly detection in Internet of Things (IoT) system logs is a critical
operational security challenge. Existing approaches rely on black-box
classifiers that provide no explanation for their predictions and fail to
correct for the severe class imbalance inherent in real-world log data. This
paper extends the work of Nguyen and Nguyen \cite{nguyen2023} by introducing
two targeted enhancements to their IoT log anomaly detection pipeline:
Synthetic Minority Oversampling Technique (SMOTE) \cite{chawla2002} to address
class imbalance, and SHapley Additive exPlanations (SHAP) \cite{lundberg2017}
to provide post-hoc interpretability. Experiments are conducted on the
Blue Gene/L (BGL) supercomputer log dataset \cite{zhu2023loghub} comprising
4{,}747{,}963 entries with a 7.34\% anomaly rate. Logs are parsed using
Drain3 \cite{he2017drain}, yielding 1{,}823 unique event templates, from
which five structural features are extracted. Four experimental conditions
evaluate Random Forest (RF) and XGBoost each trained with and without SMOTE.
Results show that SMOTE consistently improves Recall from 0.9932 to 0.9995
for both classifiers, capturing approximately 434 additional true anomalies
per 69{,}692 test anomalies at the cost of a minor Precision reduction of
0.88\%. SHAP TreeExplainer identifies template frequency as the dominant
predictive feature, revealing that logs following rare event templates are
systematically associated with anomalous behaviour --- an operationally
meaningful finding unavailable from accuracy metrics alone. Compared to the
BGL state-of-the-art benchmark established by LogAnomaly
\cite{meng2019loganomaly} (96\% F1), this work achieves 99.54\% F1 using a
simpler, non-sequential tree-based approach while additionally providing full
feature-level explainability and explicit imbalance correction.
\end{abstract}

\begin{IEEEkeywords}
IoT log anomaly detection, SMOTE, SHAP, XGBoost, Random Forest,
Drain3, BGL dataset, explainability, class imbalance
\end{IEEEkeywords}

% ────────────────────────────────────────────────────────────────────────
\section{Introduction}
\label{sec:intro}

The proliferation of Internet of Things (IoT) devices has introduced a
challenging class of security problem: the continuous, high-volume generation
of system logs that must be monitored for anomalous behaviour in near real
time. IoT devices, from industrial sensors to network infrastructure, generate
millions of log entries per day. These logs contain the earliest signals of
hardware faults, network intrusions, and system failures. Manual inspection is
infeasible at this scale, and rule-based monitoring systems fail to generalise
across the diverse and evolving log formats produced by heterogeneous IoT
deployments.

Machine learning classifiers have emerged as the dominant approach to
automated log anomaly detection. Nguyen and Nguyen \cite{nguyen2023}
demonstrated that XGBoost achieves a 99.9\% F1-score on IoT log data parsed
through the Drain algorithm \cite{he2017drain}, outperforming Random Forest,
k-Nearest Neighbours, Artificial Neural Networks, and Autoencoders on the
Thunderbird and Spirit supercomputer log datasets. However, their study leaves
two critical gaps unaddressed. First, \textit{class imbalance}: real-world log
datasets contain far fewer anomalous entries than normal ones. The authors
attribute their Autoencoder model's 29-point accuracy degradation to this
imbalance yet apply no correction. Second, \textit{explainability}: despite
strong aggregate metrics, no mechanism is provided to explain why any
individual log entry is classified as anomalous, limiting operational utility
for security practitioners who must triage and act on flagged events.

Recent work by Alam et al.\ \cite{alam2024sxad} --- the SXAD framework ---
applied Kernel SHAP to Drain-parsed HDFS log data, demonstrating XAI-driven
log anomaly detection. However, SXAD does not address class imbalance, has
not been evaluated on the BGL dataset, and uses Kernel SHAP --- a
model-agnostic approximation --- rather than the exact Shapley values
computed by TreeExplainer for tree-based models. No published study combines
SMOTE and SHAP on parsed IoT system log data, and none applies this
combination to BGL specifically.

\subsection{Problem Formulation}

Given a stream of IoT system log messages
$\mathcal{L} = \{l_1, l_2, \ldots, l_N\}$ where $l_i$ is a raw textual log
line, the binary anomaly detection task is to learn a classifier
$f: l_i \mapsto y_i \in \{0, 1\}$ where $y_i = 1$ denotes an anomalous log
entry. The challenge is twofold: (i) the class distribution is severely
imbalanced with $|\{i : y_i = 1\}| \ll |\{i : y_i = 0\}|$; and (ii) the
classifier $f$ must be interpretable, producing not only predictions but also
attributable explanations for each decision.

\subsection{Contributions}

This paper makes the following contributions:

\begin{enumerate}
  \item Full-dataset Drain3 parsing of 4{,}747{,}963 BGL log entries yielding
        1{,}823 unique event templates.
  \item SMOTE applied exclusively to the training partition to correct the
        7.34\% anomaly rate without data leakage.
  \item A controlled four-run experiment (RF and XGBoost, each with and
        without SMOTE) isolating SMOTE's effect on Recall.
  \item SHAP TreeExplainer analysis producing global feature importance,
        directional beeswarm visualisation, and instance-level waterfall
        explanations for the best-performing model.
  \item Identification of 20 specific anomalous BGL templates with associated
        operational remediation guidance, translating model output into
        actionable security intelligence.
\end{enumerate}

The remainder of this paper is organised as follows. Section~\ref{sec:related}
reviews related work. Section~\ref{sec:methodology} describes the methodology.
Section~\ref{sec:setup} presents the experimental setup.
Section~\ref{sec:results} reports and discusses results.
Section~\ref{sec:novelty} articulates the novelty of this work.
Section~\ref{sec:conclusion} concludes with future directions.

% ────────────────────────────────────────────────────────────────────────
\section{Related Work}
\label{sec:related}

\subsection{ML-Based IoT Log Anomaly Detection}

Nguyen and Nguyen \cite{nguyen2023} established the foundational pipeline
this study extends. Their approach applies Drain to parse IoT log messages
into structured event templates, extracts numeric features, and compares five
classifiers on the Thunderbird and Spirit supercomputer log datasets. XGBoost
achieves 99.9\% F1-score while the Autoencoder degrades to 70\% accuracy,
attributed to class imbalance. No SMOTE or any other imbalance correction is
applied, and no XAI method is used. This study inherits their XGBoost and RF
classifiers, hyperparameters (XGBoost: lr=0.3, max\_depth=6,
n\_estimators=100; RF: n\_estimators=20), and pipeline design to enable fair
comparison.

\textbf{Key difference:} This work adds SMOTE to correct the imbalance the
baseline acknowledges but ignores, adds SHAP to provide the explainability
the baseline completely omits, and evaluates on BGL --- a dataset not tested
in the original work.

\subsection{XAI for Log Anomaly Detection: SXAD}

Alam et al.\ \cite{alam2024sxad} proposed SXAD, which applies Kernel SHAP to
explain predictions from Decision Tree, RF, and Gradient Boosting classifiers
trained on Drain-parsed HDFS log data, reporting 99.99\% accuracy. SXAD is
the closest existing work to this study.

\textbf{Key difference:} SXAD applies no SMOTE, does not evaluate on BGL,
and uses Gradient Boosting rather than XGBoost. Critically, SXAD uses Kernel
SHAP --- a model-agnostic approximation --- whereas this work uses SHAP
TreeExplainer, which computes exact Shapley values optimised for tree-based
ensembles, producing more precise feature attributions. Furthermore, no
prior study combines SMOTE and SHAP together within a single controlled
experimental design.

\subsection{SMOTE for Imbalanced Security Data}

Talukder et al.\ \cite{talukder2024} applied SMOTE within a 10-fold
cross-validation framework to XGBoost and RF on the UNSW-NB15 network
intrusion dataset (41 continuous features), demonstrating consistent Recall
improvement on the minority attack class. Their work provides the empirical
basis for the SMOTE hypothesis in this study.

\textbf{Key difference:} UNSW-NB15 consists entirely of continuous numerical
network traffic features for which SMOTE interpolation is mathematically
well-defined. This work's features include categorical (template\_id) and
binary (has\_error) types, for which SMOTE interpolation produces values
that do not necessarily correspond to real log templates --- a limitation
this study identifies and empirically evaluates. Additionally, Talukder
et al.\ use no SHAP analysis; this study extends their finding by also
explaining which features drove the Recall improvement.

\subsection{Deep Learning for Log Anomaly Detection: DeepLog}

Du et al.\ \cite{du2017deeplog} proposed DeepLog, an LSTM-based approach
that models normal system execution as log sequences and detects deviations
using next-log prediction. DeepLog was evaluated on HDFS logs and
demonstrated strong anomaly detection capability using sequential log
patterns.

\textbf{Key difference:} DeepLog requires logs to be processed in sequential
order, which constrains its applicability to streaming IoT environments
where logs may arrive out of order or in parallel. DeepLog provides no
post-hoc explanation for individual predictions and does not address class
imbalance. This work processes each log entry independently, requires no
sequential ordering, and produces full SHAP explanations --- practical
advantages for real-world IoT deployment.

\subsection{LogAnomaly: BGL Benchmark}

Meng et al.\ \cite{meng2019loganomaly} proposed LogAnomaly, which detects
both sequential and quantitative anomalies using an LSTM with template2Vec
embeddings. LogAnomaly was evaluated on both HDFS and BGL datasets,
achieving 96\% F1-score on BGL --- establishing the published benchmark for
this dataset.

\textbf{Key difference:} LogAnomaly requires sequential log ordering and
provides no post-hoc explanation for any individual prediction. It uses
dense vector embeddings (template2Vec) rather than interpretable numerical
features, making the model's internal reasoning opaque. This work achieves
a higher F1 on BGL (99.54\% vs.\ 96\%) using a simpler, non-sequential
tree-based approach that additionally provides SHAP feature-level
explanations and explicit class imbalance correction.

\subsection{SMOTE: Foundational Oversampling Technique}

Chawla et al.\ \cite{chawla2002} introduced SMOTE, which generates synthetic
minority class samples by interpolating between existing minority instances
using k-nearest neighbours. SMOTE has since become the standard technique
for addressing class imbalance in classification tasks, with widespread
application in medical, security, and fault detection domains.

\textbf{Key difference from direct application:} The original SMOTE was
designed for continuous numerical features. This study extends SMOTE to a
domain where features include categorical template identifiers and binary
keyword indicators, documenting the interpolation limitation and evaluating
whether Recall improvement is achieved despite it --- which the empirical
results confirm.

\subsection{SHAP: Unified Model Explainability}

Lundberg and Lee \cite{lundberg2017} introduced SHAP, grounded in cooperative
game theory's Shapley values, providing consistent and locally accurate
attribution of each feature's contribution to model predictions. SHAP
TreeExplainer computes exact Shapley values for tree-based models including
XGBoost and Random Forest with polynomial time complexity.

\textbf{Key difference from prior XAI applications:} Prior applications of
SHAP in security settings (e.g., SXAD \cite{alam2024sxad}) use Kernel SHAP,
which approximates Shapley values using a model-agnostic perturbation
approach. This work uses TreeExplainer, which exploits the tree structure to
compute exact values --- more precise and computationally efficient for the
XGBoost and RF classifiers used here.

% ────────────────────────────────────────────────────────────────────────
\section{Methodology}
\label{sec:methodology}

\subsection{Dataset}

The Blue Gene/L (BGL) supercomputer log dataset is sourced from the Loghub
collection \cite{zhu2023loghub} via Zenodo (DOI: 10.5281/zenodo.8196385).
BGL contains 4{,}747{,}963 log entries collected over 214.7 days from the
Lawrence Livermore National Laboratory. Each log entry is labelled: a dash
(``-'') denotes a normal log, while any other value (e.g., APPREAD, KERDATA,
KERNELSYSERR) denotes an anomalous event. The dataset comprises 4{,}399{,}503
normal entries (92.66\%) and 348{,}460 anomalous entries (7.34\%), reflecting
severe real-world class imbalance.

Each raw log line follows the format:
\begin{verbatim}
LABEL TIMESTAMP DATE NODE DATETIME
      TYPE COMPONENT LEVEL MESSAGE
\end{verbatim}

\subsection{Log Parsing with Drain3}

Raw log messages are parsed using Drain3 \cite{he2017drain}, a Python~3
implementation of the Drain log parsing algorithm. Drain3 groups semantically
similar log messages into clusters via a fixed-depth parse tree, replacing
variable fields with wildcard placeholders (\texttt{<*>}) and assigning each
cluster a unique template identifier. Configuration parameters are set to
similarity threshold 0.4 and tree depth 4, consistent with prior work
\cite{nguyen2023}. Applied to the full BGL dataset, Drain3 discovers 1{,}823
unique event templates from 4{,}747{,}963 raw log messages.

\subsection{Feature Extraction}

Five interpretable features are derived from each parsed log entry:

\begin{enumerate}
  \item \textbf{template\_id}: The integer cluster identifier assigned by
        Drain3, encoding the structural category of the log message.
  \item \textbf{template\_length}: The number of tokens in the event
        template, capturing message structural complexity.
  \item \textbf{has\_error}: A binary indicator (0 or 1) set to 1 if the
        original message contains any of the terms \textit{error, fail,
        fatal, exception, crash} (case-insensitive).
  \item \textbf{wildcard\_count}: The count of \texttt{<*>} placeholders in
        the event template, reflecting how variable the log structure is.
  \item \textbf{template\_freq}: The occurrence count of each template\_id
        in the training partition, computed exclusively on training data and
        mapped to the test partition to prevent leakage. This feature
        encodes the rarity of each log event type.
\end{enumerate}

Formally, for a training set of size $N_{tr}$ with template assignments
$\{t_i\}_{i=1}^{N_{tr}}$, the frequency feature for template $t$ is:
\begin{equation}
  f(t) = \sum_{i=1}^{N_{tr}} \mathbb{1}[t_i = t]
  \label{eq:freq}
\end{equation}

\subsection{Train-Test Split and SMOTE Application}

The dataset is partitioned into 80\% training (3{,}798{,}370 entries) and
20\% test (949{,}593 entries) using stratified sampling with random seed
42, preserving the 7.34\% anomaly rate in both partitions. SMOTE
\cite{chawla2002} is applied exclusively to the training partition after
splitting to prevent data leakage.

For each minority class instance $\mathbf{x}_i$, SMOTE selects one of its
$k$ nearest minority-class neighbours $\mathbf{x}_{nn}$ and generates a
synthetic sample:
\begin{equation}
  \mathbf{x}_{syn} = \mathbf{x}_i + \delta \cdot
                     (\mathbf{x}_{nn} - \mathbf{x}_i),
  \quad \delta \sim U(0, 1)
  \label{eq:smote}
\end{equation}

With $k = 5$ and random seed 42, this procedure balances the training set
to 3{,}519{,}602 normal and 3{,}519{,}602 anomalous entries (50:50). The
test partition retains its original class distribution throughout all
evaluations.

\subsection{Classifiers}

Two tree-based ensemble classifiers are evaluated under two conditions
each, yielding four experimental runs:

\begin{itemize}
  \item \textbf{RF (no SMOTE)}: Random Forest, n\_estimators=20, seed=42,
        trained on imbalanced data.
  \item \textbf{RF (SMOTE)}: Random Forest, same configuration, trained on
        SMOTE-balanced data.
  \item \textbf{XGB (no SMOTE)}: XGBoost, lr=0.3, max\_depth=6,
        n\_estimators=100, seed=42, trained on imbalanced data.
  \item \textbf{XGB (SMOTE)}: XGBoost, same configuration, trained on
        SMOTE-balanced data.
\end{itemize}

Hyperparameters are held constant across conditions to isolate the effect
of SMOTE. Algorithm~\ref{alg:pipeline} formalises the complete experimental
procedure.

\begin{algorithm}[t]
\caption{Explainable IoT Log Anomaly Detection}
\label{alg:pipeline}
\begin{algorithmic}[1]
\Require Raw log file $\mathcal{L}$, similarity $\sigma$, depth $d$,
         seed $s$, neighbours $k$
\Ensure Trained model $f^*$, SHAP explanations $\Phi$
\State Initialise Drain3 parser $P$ with $\sigma, d$
\For{each line $l_i \in \mathcal{L}$}
  \State $(t_i, T_i) \gets P(\text{message}(l_i))$
        \Comment{template id, template string}
  \State $y_i \gets \mathbb{1}[\text{label}(l_i) \neq \text{`-'}]$
\EndFor
\State Extract features $\mathbf{x}_i = $ (template\_id, template\_length,
        has\_error, wildcard\_count) for all $i$
\State $(D_{tr}, D_{te}) \gets$ stratified split with seed $s$, ratio 80:20
\State Compute frequency map $f(\cdot)$ on $D_{tr}$ via Eq.~\ref{eq:freq}
\State Append template\_freq to all $\mathbf{x}_i$ via $f(t_i)$
\State $D_{tr}^{sm} \gets \text{SMOTE}(D_{tr}, k, s)$ via Eq.~\ref{eq:smote}
\For{each $(D, M) \in \{(D_{tr}, \text{RF}),(D_{tr}^{sm}, \text{RF}),
                       (D_{tr}, \text{XGB}),(D_{tr}^{sm}, \text{XGB})\}$}
  \State Train classifier $f_{D,M}$ on $D$
  \State Evaluate on $D_{te}$: Accuracy, Precision, Recall, F1
\EndFor
\State $f^* \gets \arg\max_{f_{D,M}} \text{Recall}(f_{D,M}, D_{te})$
\State $\Phi \gets$ TreeExplainer$(f^*, $ subsample of $D_{te})$
\State \Return $f^*, \Phi$
\end{algorithmic}
\end{algorithm}

Recall is the primary evaluation metric, reflecting the security-first
priority where missing a real anomaly (false negative) carries higher
operational cost than a false alarm (false positive).

\subsection{Evaluation Metrics}

Let TP, FP, TN, FN denote true positives, false positives, true negatives,
and false negatives respectively. The four standard metrics are:

\begin{align}
\text{Accuracy} &= \frac{\text{TP} + \text{TN}}
                        {\text{TP} + \text{FP} + \text{TN} + \text{FN}}
                        \label{eq:acc} \\
\text{Precision} &= \frac{\text{TP}}{\text{TP} + \text{FP}}
                         \label{eq:prec} \\
\text{Recall} &= \frac{\text{TP}}{\text{TP} + \text{FN}}
                      \label{eq:rec} \\
\text{F1} &= 2 \cdot \frac{\text{Precision} \cdot \text{Recall}}
                          {\text{Precision} + \text{Recall}}
                  \label{eq:f1}
\end{align}

\subsection{SHAP Explainability}

SHAP TreeExplainer \cite{lundberg2017} is applied to the best-performing
model on a stratified 5{,}000-entry sample of the test set. For a model
$f$, instance $\mathbf{x}$, and feature subset $S \subseteq F$ where
$F = \{1, \ldots, M\}$, the SHAP value for feature $j$ is:
\begin{equation}
  \phi_j = \sum_{S \subseteq F \setminus \{j\}}
           \frac{|S|!\,(M - |S| - 1)!}{M!}
           \bigl[v(S \cup \{j\}) - v(S)\bigr]
  \label{eq:shap}
\end{equation}
where $v(S)$ is the model output conditioned on features in $S$.
TreeExplainer computes Eq.~\ref{eq:shap} exactly in polynomial time by
exploiting tree structure, avoiding the approximation error of model-agnostic
methods such as Kernel SHAP.

Three standard visualisations are generated: (1) a global feature importance
bar chart of mean absolute SHAP values $\bar{\phi}_j = \frac{1}{n}
\sum_{i=1}^{n} |\phi_j^{(i)}|$, (2) a beeswarm plot showing the distribution
and direction of each feature's impact, and (3) a waterfall plot explaining
the prediction for a single anomalous log instance.

% ────────────────────────────────────────────────────────────────────────
\section{Experimental Setup}
\label{sec:setup}

All experiments are conducted in Google Colab using Python 3.10. The
complete software stack comprises: \texttt{pandas} and \texttt{numpy} for
data handling; \texttt{scikit-learn} for train-test splitting and
evaluation; \texttt{imbalanced-learn} for SMOTE; \texttt{XGBoost} 2.x;
\texttt{drain3} for log parsing; and \texttt{shap} for explainability.
Random seed 42 is fixed throughout all components. The BGL.log file
(approximately 708~MB) is stored in Google Drive and accessed via Drive
mount. The parsed dataset is saved as a Parquet checkpoint after Drain3
processing to enable reproducible re-runs.

Table~\ref{tab:config} summarises the full experimental configuration.

\begin{table}[h]
\caption{Experimental Configuration}
\label{tab:config}
\centering
\small
\begin{tabular}{ll}
\toprule
\textbf{Parameter} & \textbf{Value} \\
\midrule
Dataset            & BGL (Blue Gene/L) \\
Total log entries  & 4{,}747{,}963 \\
Anomaly rate       & 7.34\% (348{,}460 anomalous) \\
Log parser         & Drain3 (sim\_th=0.4, depth=4) \\
Unique templates   & 1{,}823 \\
Features           & 5 (see Methodology) \\
Train/Test split   & 80\%/20\% stratified (seed=42) \\
SMOTE k-neighbours & 5 (seed=42) \\
RF n\_estimators   & 20 \\
XGB learning rate  & 0.3 \\
XGB max depth      & 6 \\
XGB n\_estimators  & 100 \\
Primary metric     & Recall \\
SHAP method        & TreeExplainer (5{,}000-sample subset) \\
\bottomrule
\end{tabular}
\end{table}

% ────────────────────────────────────────────────────────────────────────
\section{Experimental Results and Discussion}
\label{sec:results}

\subsection{Classification Performance}

Table~\ref{tab:results} presents Accuracy, Precision, Recall, and F1-score
for all four experimental conditions evaluated on the held-out test set of
949{,}593 original BGL log entries. No synthetic samples appear in the test
set; all metrics reflect performance on real, unmodified data.

\begin{table}[h]
\caption{Classification Results on BGL Test Set (949{,}593 entries)}
\label{tab:results}
\centering
\small
\begin{tabular}{lcccc}
\toprule
\textbf{Model} & \textbf{Acc.} & \textbf{Prec.} & \textbf{Recall} & \textbf{F1} \\
\midrule
RF (no SMOTE)  & 0.9995 & 1.0000 & 0.9932 & 0.9966 \\
RF (SMOTE)     & 0.9993 & 0.9913 & 0.9995 & 0.9954 \\
XGB (no SMOTE) & 0.9995 & 1.0000 & 0.9933 & 0.9966 \\
XGB (SMOTE)*   & 0.9993 & 0.9912 & \textbf{0.9995} & 0.9954 \\
\midrule
\multicolumn{5}{l}{\footnotesize * Selected for SHAP analysis. Bold = primary metric.}\\
\bottomrule
\end{tabular}
\end{table}

Without SMOTE, both RF and XGBoost achieve Recall of 0.9932/0.9933 with
Precision of 1.0000, missing approximately 469 anomalies per 69{,}692 in
the test set. With SMOTE, Recall improves to 0.9995 for both classifiers
--- an absolute gain of 0.0063 --- while Precision decreases marginally to
0.9912/0.9913. In absolute terms, SMOTE-trained models catch approximately
434 additional true anomalies at the cost of approximately 613 extra false
positives.

These results confirm the hypothesis grounded in Talukder et al.\
\cite{talukder2024}: SMOTE consistently improves Recall for tree-based
ensemble classifiers on imbalanced security log data. The identical response
of RF and XGBoost to SMOTE indicates the effect is driven by data balancing
rather than model-specific characteristics --- a finding that would not be
observable from a single-classifier experiment. The trade-off --- 0.63\%
Recall gain for 0.88\% Precision reduction --- is operationally favourable
in security contexts, where a missed anomaly carries higher cost than a
false alarm.

\subsection{Comparison with Published Benchmarks}

Table~\ref{tab:benchmark} places the results in context with published
benchmarks.

\begin{table}[h]
\caption{Benchmark Comparison}
\label{tab:benchmark}
\centering
\small
\begin{tabular}{llcc}
\toprule
\textbf{Method} & \textbf{Dataset} & \textbf{F1} & \textbf{Source} \\
\midrule
LogAnomaly \cite{meng2019loganomaly} & BGL & 0.9600 & IJCAI 2019 \\
Nguyen \& Nguyen \cite{nguyen2023}   & Thunderbird & 0.9990 & CNIOT 2023 \\
SXAD \cite{alam2024sxad}             & HDFS & N/A\textsuperscript{a} & IEEE Access 2024 \\
\textbf{XGB+SMOTE (this work)}       & \textbf{BGL} & \textbf{0.9954} & \textbf{This paper} \\
\bottomrule
\multicolumn{4}{l}{\footnotesize \textsuperscript{a}SXAD reports 99.99\% Accuracy; F1 not reported.}\\
\end{tabular}
\end{table}

Compared to LogAnomaly \cite{meng2019loganomaly}, which reports 96\% F1 on
BGL using an LSTM-based sequential approach, this work achieves 99.54\%
F1-score on the same dataset using a simpler, non-sequential tree-based
approach while additionally providing SHAP interpretability and explicit
imbalance correction. The 3.54 percentage-point improvement is achieved
without sequential log ordering and with full feature-level explainability
--- a combination not available in any prior BGL study.

\subsection{SHAP Feature Importance}

SHAP TreeExplainer is applied to XGB (SMOTE) on a stratified 5{,}000-entry
sample of the test set. The global feature importance ranking from highest
to lowest mean absolute SHAP value is: \textbf{template\_freq},
\textbf{template\_id}, \textbf{template\_length}, \textbf{wildcard\_count},
\textbf{has\_error}. Fig.~\ref{fig:shap_importance} shows the global
importance bar chart.

\begin{figure}[h]
  \centering
  \includegraphics[width=\columnwidth]{shap_importance.png}
  \caption{SHAP global feature importance for XGBoost trained with SMOTE.
  Mean absolute SHAP values across 5{,}000 stratified test samples.}
  \label{fig:shap_importance}
\end{figure}

Template frequency is the dominant predictive feature, followed by template
identifier. Together, these two structural features account for the majority
of the model's predictive signal. The error keyword indicator (has\_error)
--- the most intuitive anomaly signal --- ranks last, contributing least to
predictions.

The beeswarm plot in Fig.~\ref{fig:shap_beeswarm} reveals a clear
directional pattern for template\_freq: low frequency values (blue dots)
spread far to the right of the zero SHAP axis, strongly pushing predictions
toward anomaly; high frequency values (red dots) cluster to the left,
pushing toward normal. This reveals the model's core detection logic:
\textit{logs following rare event templates are systematically anomalous,
while high-frequency templates reflect routine system behaviour}.

\begin{figure}[h]
  \centering
  \includegraphics[width=\columnwidth]{shap_beeswarm.png}
  \caption{SHAP beeswarm plot. Low template\_freq values (blue) push
  predictions toward anomaly (right); high values (red) push toward normal.}
  \label{fig:shap_beeswarm}
\end{figure}

This finding carries direct operational value. A security practitioner can
use template frequency as a prioritisation signal: log entries from rare
templates warrant immediate investigation, while high-frequency template
entries can be deprioritised. This interpretable rule cannot be derived from
the black-box predictions of the baseline \cite{nguyen2023} or from
LogAnomaly \cite{meng2019loganomaly}.

Critically, has\_error --- the most intuitive anomaly signal --- ranks
last. Explicit error keyword matching contributes less predictive signal
than structural template features, challenging the assumption underlying
most rule-based log monitoring systems that lexical indicators are the
primary discriminator.

\subsection{Instance-Level Explainability}

Fig.~\ref{fig:shap_waterfall} presents a SHAP waterfall plot for a single
anomalous log entry from the test set. The waterfall decomposes the model
output from the base value $E[f(\mathbf{x})]$ to the final anomaly
probability, showing each feature's additive contribution. For the selected
instance, the low template\_freq value is the largest contributor pushing
the prediction toward anomaly, followed by template\_id. This instance-level
explanation allows a security analyst to understand precisely which
structural characteristics triggered each alert --- a capability absent from
all prior BGL anomaly detection work.

\begin{figure}[h]
  \centering
  \includegraphics[width=\columnwidth]{shap_waterfall.png}
  \caption{SHAP waterfall plot for a single anomalous BGL log entry, showing
  each feature's additive contribution to the model output.}
  \label{fig:shap_waterfall}
\end{figure}

\subsection{Template-Level Explainability}

Combining SHAP rankings with Drain3 template frequency analysis reveals 20
specific anomalous templates. Table~\ref{tab:templates} presents the top
occurrences.

\begin{table}[h]
\caption{Top Anomalous Templates in BGL Test Set}
\label{tab:templates}
\centering
\small
\begin{tabular}{clcc}
\toprule
\textbf{T-ID} & \textbf{Template (truncated)} & \textbf{Hits} & \textbf{Freq} \\
\midrule
18   & data TLB error interrupt               & 142 & 122{,}358 \\
192  & data storage interrupt                 & 79  & 50{,}781 \\
587  & Lustre mount FAILED                    & 37  & 21{,}085 \\
1108 & ciod: Error reading...Link severed     & 29  & 23{,}050 \\
568  & rts: kernel terminated reason 1001     & 24  & 14{,}876 \\
449  & ciod: Error after LOGIN\_MESSAGE       & 21  & 18{,}568 \\
17   & ciod: failed to read control stream    & 11  & 4{,}792 \\
565  & ciodb exited abnormally: Aborted       & 1   & 18 \\
\bottomrule
\end{tabular}
\end{table}

The rarest templates (frequency 1--4) receive the highest SHAP anomaly
scores consistent with the template\_freq finding. For example, Template
565 (\textit{ciodb exited abnormally due to signal: Aborted}) appears only
18 times in training yet is flagged with high anomaly confidence --- a
daemon crash event warranting immediate escalation. Template 824
(\textit{ddr: redundant bit steering failed, sequencer timeout}), with
frequency 1, indicates a DDR memory redundancy failure with data integrity
implications.

\subsection{Threats to Validity}

Following standard guidance for empirical software engineering research, we
discuss four classes of threats to validity.

\textbf{Internal validity.} All experiments use a single random seed (42).
While this ensures reproducibility, it does not establish statistical
significance across multiple runs. Future work should report mean and
confidence intervals across multiple seeds.

\textbf{External validity.} Results are obtained on a single dataset (BGL),
which is a supercomputer log corpus rather than a heterogeneous IoT device
log set. Generalisation to consumer IoT devices, industrial sensors, or
network edge logs remains to be empirically validated.

\textbf{Construct validity.} The five extracted features are structural and
frequency-based; semantic content beyond keyword matching is not captured.
SMOTE interpolates between template\_id values --- a categorical feature
--- producing synthetic samples that do not necessarily map to real log
templates. While empirical Recall improvement is observed despite this,
the synthetic samples are not semantically meaningful as standalone log
entries.

\textbf{Conclusion validity.} The SHAP analysis is computed on a 5{,}000-
entry stratified subset of the test set due to computational constraints
in the Colab environment. A full-test-set SHAP evaluation may yield
marginally different feature importance rankings.

\subsection{Reproducibility}

To support reproducibility, this work fixes random seed 42 throughout, uses
publicly available datasets and libraries, and reports all hyperparameters
in Table~\ref{tab:config}. The complete experimental procedure is given in
Algorithm~\ref{alg:pipeline}. The BGL dataset is available at
\url{https://zenodo.org/records/8196385}. All software used (pandas,
scikit-learn, imbalanced-learn, XGBoost, drain3, shap) is open source.
The complete source code, parsed dataset checkpoint, and SHAP figures are
publicly available at
\url{https://github.com/syedfaisal2407/iot-anomaly-detection-bgl}.

% ────────────────────────────────────────────────────────────────────────
\section{Novelty}
\label{sec:novelty}

This work makes four novel contributions not present in any single prior
study:

\textbf{(N1) First SMOTE+SHAP combination on IoT log data.} SXAD
\cite{alam2024sxad} applies SHAP to log data without SMOTE. Talukder et al.\
\cite{talukder2024} apply SMOTE without SHAP. No prior study combines both
within a controlled experiment, as this work does.

\textbf{(N2) First evaluation of SMOTE+SHAP on BGL.} The BGL dataset has
been benchmarked by LogAnomaly \cite{meng2019loganomaly} and studied by
He et al.\ \cite{he2017drain} for parsing evaluation, but no prior study
applies SMOTE or SHAP to BGL specifically.

\textbf{(N3) Exact SHAP values via TreeExplainer.} Prior log anomaly XAI
work \cite{alam2024sxad} uses Kernel SHAP approximations. This work uses
SHAP TreeExplainer, which computes exact Shapley values for XGBoost and
RF, producing more reliable feature attributions.

\textbf{(N4) Template-level operational explainability.} By combining SHAP
rankings with Drain3 frequency analysis, this work identifies 20 specific
anomalous BGL templates with associated system error explanations and
remediation recommendations --- translating model output into actionable
security intelligence. No prior BGL anomaly detection work provides this
level of operational interpretability.

% ────────────────────────────────────────────────────────────────────────
\section{Conclusion}
\label{sec:conclusion}

This paper extended the IoT log anomaly detection pipeline of Nguyen and
Nguyen \cite{nguyen2023} by integrating SMOTE \cite{chawla2002} for class
imbalance correction and SHAP TreeExplainer \cite{lundberg2017} for
post-hoc interpretability, evaluated on the entire 4{,}747{,}963-entry BGL
dataset. The controlled four-run experiment demonstrated that SMOTE
consistently improves Recall from 0.9932 to 0.9995 for both RF and XGBoost,
catching approximately 434 additional anomalies per test set at the cost of
613 additional false positives. XGBoost trained with SMOTE achieved the
best Recall (0.9995) and was selected for SHAP analysis.

SHAP TreeExplainer identified template frequency as the dominant predictive
feature, establishing that rare event templates are the primary anomaly
signal --- a finding with direct operational value for security
practitioners. Compared to the BGL state-of-the-art established by
LogAnomaly \cite{meng2019loganomaly} (96\% F1), this work achieves 99.54\%
F1 with a simpler non-sequential architecture while additionally providing
full feature-level explainability and explicit imbalance correction.

Future work should address SMOTE compatibility with categorical features
(e.g., SMOTENC), extend SHAP analysis to the full test set, validate
findings on additional IoT log datasets beyond BGL, perform multi-seed
statistical significance testing, and explore multi-class anomaly
classification to automatically recommend remediation actions.

% ────────────────────────────────────────────────────────────────────────
\section*{Ethics Statement}

This research uses the publicly available BGL dataset
\cite{zhu2023loghub}, which contains system logs from a Lawrence Livermore
National Laboratory supercomputer. The dataset contains no personally
identifiable information. All software used is open source. The research
is conducted for academic purposes and to advance the state of practice in
operational security; no offensive applications are intended.

\section*{Acknowledgment}

The author thanks the Department of Computer Information Systems and
Cybersecurity at Auburn University at Montgomery for supporting this
research.

% ────────────────────────────────────────────────────────────────────────
\begin{thebibliography}{00}

\bibitem{nguyen2023}
H.-T. Nguyen and V. Q. Nguyen,
``Anomaly Detection in IoT System Logs Using Machine Learning and Deep
Learning Approaches,''
in \textit{Proc. ACM Int. Conf. on Computing and Network Communications
(CNIOT)}, Xiamen, China, 2023, pp. 969--974.
doi: 10.1145/3603781.3604223.

\bibitem{alam2024sxad}
S. Alam, K. Kifayat, C. Sampedro, V. Karovi\v{c}, and H. Naeem,
``SXAD: Shapley eXplainable AI-Based Anomaly Detection Using Log Data,''
\textit{IEEE Access}, vol.~12, pp. 95659--95672, 2024.
doi: 10.1109/ACCESS.2024.3425472.

\bibitem{talukder2024}
Md. A. Talukder \textit{et al.},
``Machine Learning-Based Network Intrusion Detection for Big and Imbalanced
Data Using Oversampling and Learning Algorithm,''
\textit{Journal of Big Data}, vol.~11, no.~1, 2024.
doi: 10.1186/s40537-024-00886-w.

\bibitem{meng2019loganomaly}
W. Meng \textit{et al.},
``LogAnomaly: Unsupervised Detection of Sequential and Quantitative
Anomalies in Unstructured Logs,''
in \textit{Proc. IJCAI}, Macao, 2019, pp. 4739--4745.
doi: 10.24963/ijcai.2019/658.

\bibitem{du2017deeplog}
M. Du, F. Li, G. Zheng, and V. Srikumar,
``DeepLog: Anomaly Detection and Diagnosis from System Logs Through Deep
Learning,''
in \textit{Proc. ACM CCS}, Dallas, TX, 2017, pp. 1285--1298.
doi: 10.1145/3133956.3134015.

\bibitem{chawla2002}
N. V. Chawla, K. W. Bowyer, L. O. Hall, and W. P. Kegelmeyer,
``SMOTE: Synthetic Minority Over-sampling Technique,''
\textit{Journal of Artificial Intelligence Research}, vol.~16, pp. 321--357,
2002.
doi: 10.1613/jair.953.

\bibitem{lundberg2017}
S. M. Lundberg and S.-I. Lee,
``A Unified Approach to Interpreting Model Predictions,''
in \textit{Proc. NeurIPS}, Long Beach, CA, 2017, pp. 4765--4774.

\bibitem{he2017drain}
P. He, J. Zhu, Z. Zheng, and M. R. Lyu,
``Drain: An Online Log Parsing Approach with Fixed Depth Tree,''
in \textit{Proc. IEEE ICWS}, Honolulu, HI, 2017, pp. 33--40.
doi: 10.1109/ICWS.2017.8029742.

\bibitem{zhu2023loghub}
J. Zhu, S. He, P. He, J. Liu, and M. R. Lyu,
``Loghub: A Large Collection of System Log Datasets for AI-Driven Log
Analytics,''
in \textit{Proc. IEEE ISSRE}, 2023.
doi: 10.1109/ISSRE5003.2020.00012.

\bibitem{breiman2001random}
L. Breiman,
``Random Forests,''
\textit{Machine Learning}, vol.~45, no.~1, pp. 5--32, 2001.
doi: 10.1023/A:1010933404324.

\bibitem{chen2016xgboost}
T. Chen and C. Guestrin,
``XGBoost: A Scalable Tree Boosting System,''
in \textit{Proc. ACM SIGKDD}, San Francisco, CA, 2016, pp. 785--794.
doi: 10.1145/2939672.2939785.

\end{thebibliography}

\balance

\end{document}
